\section{Alignments and tagged rows}\label{sec:preliminaries}

We first relate edit distance to increasing subsequence matchings, then
transfer a separation estimate for tagged child words to an estimate for
rows under arbitrary alignments. This is the only general tool needed
before constructing the circle words.

\subsection{Subsequence deficit and tagged rows}
A common subsequence is a collection of equal-symbol pairs whose position
indices increase strictly on both sides. Write $\LCS(x,y)$ for its maximum
size. For equal-length words of length $n$, put
\[
 r_n(x,y)=n-\LCS(x,y).
\]
This is the number of unmatched positions on either side of a longest
common subsequence. For arbitrary lengths put
$q(x,y)=\max(|x|,|y|)-\LCS(x,y)$.

\begin{lemma}[Metric comparisons]\label{pre:metric}
For all finite words $x,y$,
\[
 q(x,y)\leq\ED(x,y)\leq2q(x,y),\qquad
 \ED(x,y)\leq |x|+|y|-2\LCS(x,y)\leq2\ED(x,y).
\]
The middle expression in the second comparison is the insertion--deletion
distance. In particular, $r_n\leq\ED\leq2r_n$ at equal length.
\end{lemma}
\begin{proof}
A script of $k$ edits leaves at least $|x|-k$ original symbols unchanged
and in order. Apply this observation to the script and its reverse to get
$q\leq\ED$. Deleting outside a longest common subsequence and inserting
the missing symbols costs $|x|+|y|-2\LCS(x,y)\leq2q$.
The survivor count shows that no insertion--deletion script can be shorter.
Replacing every substitution by one deletion and one insertion proves the
remaining factor-two comparison. These are the trace identities of
Wagner and Fischer \cite{WagnerFischer1974}, with the metric convention
made explicit here.
\end{proof}

Here is the alignment fact used by both constructions. A row consists of
$t$ child blocks of equal length, in the same fixed child order in both
words. Every letter carries its child tag, so a matched pair must preserve
the child index. A source row is \emph{split} if its matched positions enter
at least two target rows.

\begin{lemma}[Charging split tagged rows]\label{pre:tagged-rows}
Let $R,L\geq1$ and $t\geq2$ be integers, and let
$\Gamma_1,\ldots,\Gamma_t$ be pairwise disjoint alphabets. Each of two words
consists of $R$ rows, each containing in order a block in $\Gamma_i^L$ for
$i=1,\ldots,t$. Thus every matched pair preserves its child index. The
common length is $N=RtL$. Fix an increasing matching with $E$ unmatched positions
on each side. Suppose every nonsplit source row loses at least $AtL$ source
positions, where $0\leq A\leq1$; rows with no matches are included in this
hypothesis. Then
\begin{equation}\label{pre:row-bound}
 E\geq A\left(N-\frac{tE}{t-1}\right),\qquad
 E\geq\frac{AN}{1+tA/(t-1)}.
\end{equation}
\end{lemma}
\begin{proof}
In a split source row, choose consecutive matched positions at a transition
between target rows. They are consecutive in the whole matching because all
source positions between them remain in this same row. Let their child
indices be $i\leq j$ and their target row indices be $a<a'$. Tags force the
same $i,j$ in the target. The open target interval between the two matches
contains at least
\[
 \bigl((a'-a)t+j-i-1\bigr)L\geq(t-1)L
\]
positions: within the endpoint blocks, the smallest possible gap occurs
when the first match is last in its block and the second is first in its
block. Every position of this interval is unmatched, by consecutiveness.
The target intervals selected from different source rows are disjoint, by
the order of the matching. There are at most $E/((t-1)L)$ split rows, of
total source length at most $tE/(t-1)$.

Every remaining row loses at least the fraction $A$ of its own source
positions. These losses are disjoint even when several source rows use one
target row. Summing them gives the first inequality in~\eqref{pre:row-bound};
rearrangement gives the second.
\end{proof}

\begin{figure}[ht]
\centering
\begin{tikzpicture}[x=.72cm,y=.75cm,>=Stealth,font=\small]
 \node[anchor=east] at (-.25,2) {source row};
 \node[anchor=east] at (-.25,0) {target rows};
 \foreach \i in {0,...,4} {
   \draw[fill=blue!5] (\i,1.65) rectangle ++(1,.7);
   \pgfmathtruncatemacro{\tag}{\i+1}
   \node at (\i+.5,2) {$\tag$};
 }
 \foreach \i in {0,...,9} {
   \draw[fill=black!3] (\i,-.35) rectangle ++(1,.7);
   \pgfmathtruncatemacro{\tag}{mod(\i,5)+1}
   \node at (\i+.5,0) {$\tag$};
 }
 \draw[very thick] (5,-.5)--(5,.5);
 \draw[->,thick,blue!60!black] (.5,1.6)--(.5,.4);
 \draw[->,thick,blue!60!black] (3.5,1.6)--(8.5,.4);
 \draw[<->] (.5,-.8)--(8.5,-.8);
 \node[fill=white,inner sep=2pt] at (4.5,-.8)
   {unmatched gap, length at least $(t-1)L$};
 \node at (2.5,-1.4) {one row};
 \node at (7.5,-1.4) {next row};
\end{tikzpicture}

\caption{Two consecutive matches from one source row enter different target
rows. The labels are child tags. The open target interval between the arrows
is unmatched and contains at least $(t-1)L$ positions. The drawing shows
$t=5$ schematically; only the two relevant matches are shown.}
\label{fig:split-row}
\end{figure}
